% HD-F1-T14. GPT-6 Astra Pro, /root/transport; 2026-11-22 UTC.
\section[HD-F1-T14: nine Lie-type source exclusions]{HD-F1-T14: nine Lie-type families excluded at the original source}
\label{sec:f1-low-lie-rank-source-exclusion}

The remaining T2-OTHER source gate has no almost simple
projective coefficient image with socle in any row of the
following table. The rank bound is on the original complex
coefficient representation, of dimension \(m\); it is not
the defining-characteristic natural-module dimension used
to name the finite simple group.
\[
\begin{array}{c|c|c|c}
 S&\text{parameters}&\alpha_S(x)\le a&m\le\frac{6a}{12-2a}\\ \hline
 PSL_2(q)&q\ge4&5&15\\
 PSL_3(q)&q\ge2&4&6\\
 PSU_3(q)&q\ge3&4&6\\
 PSL_5(q)&q\ge2&5&15\\
 PSU_5(q)&q\ge2&5&15\\
 PSp_4(q)&q\ge4&5&15\\
 G_2(q)&q\ge3&5&15\\
 {}^2B_2(q)&q=2^{2e+1},\ e\ge1&3&3\\
 {}^2G_2(q)&q=3^{2e+1},\ e\ge1&3&3
\end{array}
 \tag{F1T14.1}
\]
All displayed parameters give simple groups. In particular,
the \(PSp_4(3)\) exception, the dimension-four linear and
unitary groups, and the higher exceptional groups are not
included by this statement. The established reduction
applies to every central lift and every inertia order over
each of the eight remaining coefficient fields.

\paragraph{1. Exact published bounds and the passage from prime to arbitrary order.}
Use Guralnick--Saxl, \emph{Generation of finite almost simple
groups by conjugates}, \emph{J. Algebra}~268 (2003),
519--571, DOI \texttt{10.1016/S0021-8693(03)00182-0}.
The exact inputs, with their exceptions retained, are:
\begin{itemize}
\item Lemma~3.1, pp.~527--529: the uniform bound five for
\(PSL_2(q)\), including the field-involution exception
at \(q=9\).
\item Lemmas~3.2 and~3.3, pp.~529--531: the bounds four
for \(PSL_3(q)\) and \(PSU_3(q)\), including graph-field
and \(PSU_3(3)\) involutions.
\item Theorem~4.2, p.~533: the bound five for the two
dimension-five linear and unitary families. The symplectic
transvection exception of that theorem is absent in these
two families.
\item Theorem~4.3, p.~533: for \(PSp_4(q)\), the ordinary
bound is \(2+3=5\), and its transvection bounds are four
or five. Its separate \(q=3\) bound six is excluded from
the table.
\item Theorem~5.1 and Proposition~5.6, pp.~541,545--546:
\(G_2(q)\) has untwisted rank two and bound five. The
stated \(F_4\) exception has no application here.
\item Proposition~5.8, p.~547: the two displayed Suzuki
and Ree families have bound three.
\end{itemize}
The relevant full proofs were read in the primary article:
Lemmas~3.1--3.4, the linear, unitary and symplectic proof
of Theorems~4.1--4.3 through p.~539, and the complete proofs
of Propositions~5.6 and~5.8. The source and comparison
record specifies this scope, without claiming an independent
proof of all finite-group classification inputs.

For completeness, the bounds stated first for prime-order
automorphisms give the same bounds for every nonidentity
\(x\in\operatorname{Aut}(S)\). Choose a power
\(y=x^t\ne1\) of prime order. Conjugates
\(y^{s_1},\ldots,y^{s_a}\) that generate a group containing
\(S\) lie in the subgroup generated by the actual
\(x^{s_1},\ldots,x^{s_a}\), because
\[
 y^{s_i}=(x^{s_i})^t.
 \tag{F1T14.2}
\]
Thus \(\alpha_S(x)\le\alpha_S(y)\le a\). The conjugating
elements lie in \(\langle S,x\rangle\), which is a subgroup
of the actual projective coefficient image. No change to
the original branch tuple is involved in this argument.

\paragraph{2. Apply the bound inside the original coefficient space.}
Let \(G\subset GL_K(V)\) be the actual faithful absolutely
irreducible coefficient image, let \(Z=Z(G)\) be its scalar
centre, and let \(S\) be the socle of its almost simple
projective image. For the inverse image \(P\) of \(S\),
the exact equalities
\[
 P=ZL,\quad L=P'=[L,L],\quad V_{\mathbf C}^L=0
 \tag{F1T14.3}
\]
and the actual conjugate-lift argument in T13, paragraph~1,
apply to every simple \(S\) in (F1T14.1). Specifically,
the group generated by lifted conjugates contains \(L\),
and the intersection of their equal-eigenvalue spaces is
an \(L\)-fixed space. Therefore every eigenspace of every
noncentral actual \(g\in G\) has codimension at least
\(m/a\). A nonidentity scalar has moving rank \(m\),
which satisfies the same inequality. This argument does
not require any hypothesis on the Schur multiplier of
\(S\), nor irreducibility after restriction to \(L\).

For the original connected full sphere cover, denote its
true inertia matrices by \(g_j\). Its original parabolic
dimension six and period rank nine give precisely
\[
 2m+6=\sum_{j=1}^b\operatorname{rank}(g_j-1),\qquad
 b\ge12.
 \tag{F1T14.4}
\]
Consequently
\[
 2m+6\ge bm/a\ge12m/a,
 \qquad m\le\frac{6a}{12-2a}.
 \tag{F1T14.5}
\]
These are the last-column bounds in (F1T14.1), all strictly
smaller than the remaining coefficient lower bound
\(m\ge27\). This proves each of the nine exclusions.

\paragraph{3. The exact geometric and cohomological maps.}
The map of covers used here is the identity of the original
\(C\to\mathbf P^1\). The map of middle-extension coefficient
sheaves is the identity; its branch stalks are
\(V^{g_j}\), whose dimensions are bounded in that same
\(V\) by (F1T14.3). The finite Hurwitz map to the pointed
sphere base has dimension \(b-3\), giving the second
inequality in (F1T14.4). The source has not been replaced
by a representation in the finite group's defining
characteristic.

The prescribed rational primitive idempotent in
\(\mathbf Q[G]\) and its actual integer-cleared Jacobian
image remain the same. The induced maps on the factor,
restricted theta polarization, Poincare biextension,
parabolic cup pairing, Hodge filtration and period
differential are identity maps with multiplier one.
The same identity transports its Grothendieck and Adams
realizations. Thus these are exclusions of actual source
families for the original sixfold corner, and do not count
as additional algebraic detector components.
